\documentclass[10pt,a4paper]{amsart}
\usepackage[margin=19mm]{geometry}
\usepackage{amsmath,amssymb,mathtools}
\usepackage[scr=rsfs,cal=euler]{mathalfa}
\usepackage{xfrac}
\usepackage[unicode,hidelinks,bookmarks=false]{hyperref}
\newcommand{\ie}{i.e.\ }
\newcommand{\cf}{cf.\ }
\newcommand{\resp}{respectively\ }
\newcommand{\Subsection}{\subsection}
\providecommand{\nobreakdash}{\nobreak\hbox{-}}
\setlength{\emergencystretch}{2em}
\multlinegap0pt
\usepackage{bm}
\usepackage{bbm}
\usepackage{dsfont}
\usepackage{xspace}
\usepackage{cancel}
\usepackage{mathtools}
\usepackage{amsthm}
\makeatletter
\@ifundefined{theorem}{\newtheorem{theorem}{Theorem}}{}
\@ifundefined{lemma}{\newtheorem{lemma}[theorem]{Lemma}}{}
\makeatother
\DeclareMathOperator{\diam}{diam}
\title[Counterexamples to the Corsten-Frankl conjecture on diameter]{Counterexamples to the Corsten-Frankl conjecture on diameter-Ramsey simplices}
\author{Yaping Mao}
\date{Source: arXiv:2604.19126v1 (CC BY 4.0)}
\begin{document}
\maketitle

\noindent\emph{Source and licence.} Counterexamples to the Corsten-Frankl conjecture on diameter-Ramsey simplices. Version arXiv:2604.19126v1 (CC BY 4.0), distributed under the Creative Commons Attribution 4.0 International licence, as recorded in the arXiv licence metadata for this exact version and shown on the abstract page https://arxiv.org/abs/2604.19126v1. Full rights notice: licenses/arxiv-2604.19126-CC-BY-4.0.txt.\par\smallskip

\noindent\textit{Editorial scope.} Counted unit: the Theorem whose source label is \texttt{thm:criterion} in the article (source0.tex lines 167--188, source label \texttt{thm:criterion}), delivered together with the article's own complete original proof of it (source0.tex lines 190--258, one \texttt{proof} environment of 69 lines). No mathematics is written, repaired or re-derived: statement and proof are the article's own text line for line. Required context retained uncounted: lem:regular (lines 109--111); lem:product (lines 117--126). Every definition, display and label of those lines is retained unchanged. Omitted from this package and not redistributed: 1--108, 112--116, 127--166, 189--189, 259--520. Nothing inside a retained excerpt is omitted or altered: every retained body is delivered exactly as the article's lines are written, so no omission receipt of any kind accompanies this package. Citations: the retained excerpts carry no citation command at all, so no bibliography entry of the article is transcribed and no reference is redistributed. Reference adaptation: none was needed and none was made; every reference inside the delivered unit resolves inside it, so no unresolved reference or citation is produced. Printed slips kept literal: no printed word, symbol, hypothesis or number of the counted statement or of its proof was altered. Layout and class: the article's own class is \texttt{elsarticle} with packages including fontenc, inputenc, lmodern, microtype, amsmath, amssymb, amsthm, mathtools, hyperref; the wrapper is the lane's pinned \texttt{amsart} editorial wrapper at 10pt on A4, which restates the theorem-like environments the retained text needs and adds the article's own macro definitions that the retained text needs (\texttt{\textbackslash diam}), with extra wrapper packages (bm, bbm, dsfont, xspace, cancel, mathtools). The delivered numbering is the unit's own. Assets: no retained span contains a figure environment, a graphics-inclusion command, a PDF-page inclusion, a table or a TikZ picture; no image file is copied into this package, the package declares no asset dependency and no image file is redistributed with it. Licence: version arXiv:2604.19126v1 of this article is distributed under the Creative Commons Attribution 4.0 International licence, as recorded in the arXiv licence metadata for this exact version and shown on the abstract page https://arxiv.org/abs/2604.19126v1. A full rights notice is delivered at licenses/arxiv-2604.19126-CC-BY-4.0.txt. Characters: every retained excerpt is pure ASCII, so the delivered engine, XeLaTeX with its default Latin Modern text font, reports no missing character.
\begin{lemma}\label{lem:regular}
Every regular simplex is diameter-Ramsey.
\end{lemma}

\begin{lemma}\label{lem:product}
If $X\subset\mathbb R^m$ and $Y\subset\mathbb R^n$ are diameter-Ramsey, then their Cartesian product
\[
X\times Y:=\{(x,y):x\in X,\ y\in Y\}\subset\mathbb R^{m+n},
\]
with the Euclidean product metric, is diameter-Ramsey. Moreover,
\[
\diam(X\times Y)^2=\diam(X)^2+\diam(Y)^2.
\]
\end{lemma}

\begin{theorem}\label{thm:criterion}
Let $A=\{p_1,\dots,p_n\}$ be an $(n-1)$-simplex with $n\ge 2$, and let
\[
D:=\diam(A).
\]
Relabel the vertices so that $\|p_1-p_2\|=D$. Let
\[
\mathcal B:=\{B\subseteq[n]: |B|\ge 2\ \text{and}\ \{1,2\}\nsubseteq B\}.
\]
Suppose that there are nonnegative real numbers $\alpha_B$, indexed by $B\in\mathcal B$, such that
\begin{equation}\label{eq:deficits}
D^2-\|p_i-p_j\|^2
=
\sum_{\substack{B\in\mathcal B\\ \{i,j\}\subseteq B}}\alpha_B
\qquad (1\le i<j\le n),
\end{equation}
and
\begin{equation}\label{eq:mass}
\sum_{B\in\mathcal B}\alpha_B\le D^2.
\end{equation}
Then $A$ is diameter-Ramsey.
\end{theorem}

\begin{proof}
Define the reserve mass
\[
\alpha_0:=D^2-\sum_{B\in\mathcal B}\alpha_B\ge 0.
\]
We shall realize $A$ inside a finite Cartesian product of regular simplices whose diameter is $D$.

If $\alpha_0>0$, let $S_0$ be a regular $(n-1)$-simplex of side length $\sqrt{\alpha_0}$ with distinct vertices
\[
v_1^{(0)},\dots,v_n^{(0)}.
\]
If $\alpha_0=0$, let $S_0$ be a singleton and put $v_1^{(0)}=\cdots=v_n^{(0)}$ equal to its unique point.

For each $B\in\mathcal B$ with $\alpha_B>0$, let $S_B$ be a regular simplex of side length $\sqrt{\alpha_B}$ with vertex set
\[
\{u_B\}\cup\{v_i^{(B)}:i\notin B\}.
\]
This is possible because the displayed set has $n-|B|+1$ vertices. Define
\[
q_i^{(B)}:=
\begin{cases}
 u_B, & i\in B,\\[1mm]
 v_i^{(B)}, & i\notin B,
\end{cases}
\qquad (i\in[n]),
\]
and set $q_i^{(0)}:=v_i^{(0)}$.

Now put
\[
\mathcal R:=S_0\times \prod_{\substack{B\in\mathcal B\\ \alpha_B>0}} S_B
\]
and
\[
q_i:=\bigl(q_i^{(0)},(q_i^{(B)})_B\bigr)\in\mathcal R
\qquad (i\in[n]).
\]

Fix $1\le i<j\le n$. In the factor $S_B$, the points $q_i^{(B)}$ and $q_j^{(B)}$ coincide if and only if $\{i,j\}\subseteq B$; otherwise they are distinct vertices of a regular simplex and are at squared distance $\alpha_B$. Therefore
\begin{align*}
\|q_i-q_j\|^2
&=\alpha_0+
\sum_{\substack{B\in\mathcal B\\ \alpha_B>0\\ \{i,j\}\nsubseteq B}}\alpha_B \\
&=D^2-
\sum_{\substack{B\in\mathcal B\\ \{i,j\}\subseteq B}}\alpha_B \\
&=\|p_i-p_j\|^2,
\end{align*}
where the last equality is \eqref{eq:deficits}. Thus $\{q_1,\dots,q_n\}$ is congruent to $A$.

Moreover, no set $B\in\mathcal B$ contains both $1$ and $2$. Hence $q_1$ and $q_2$ are distinct in every nontrivial factor, and
\[
\|q_1-q_2\|^2
=\alpha_0+
\sum_{B\in\mathcal B}\alpha_B
=D^2.
\]
The same sum is the squared diameter of the product $\mathcal R$, so
\[
\diam(\mathcal R)=D.
\]

By Lemma~\ref{lem:regular}, every nontrivial factor of $\mathcal R$ is diameter-Ramsey, and a singleton is trivially diameter-Ramsey. Repeated use of Lemma~\ref{lem:product} shows that $\mathcal R$ is diameter-Ramsey. Therefore, for every positive integer $q$, there exists a finite set $W_q$ such that
\[
\diam(W_q)=\diam(\mathcal R)=D
\qquad\text{and}\qquad
W_q\to(\mathcal R)_q.
\]
Since $\mathcal R$ contains the subset $\{q_1,\dots,q_n\}$ congruent to $A$, every monochromatic copy of $\mathcal R$ contains a monochromatic copy of $A$. Hence $W_q\to(A)_q$ for every $q$, and $A$ is diameter-Ramsey.
\end{proof}

\end{document}
